\documentclass[12pt,a4paper]{article}
\usepackage{amsmath,amssymb}
\usepackage{graphicx}
\usepackage{geometry}
\usepackage{float}
\usepackage{booktabs}
\usepackage{xcolor}
\usepackage{hyperref}


\usepackage{graphicx}
\usepackage{float}   

\usepackage{xepersian}
\settextfont{B-NAZANIN.TTF} 
\setlatintextfont[Scale=0.85]{Times New Roman}

\renewcommand{\labelitemi}{$\bullet$}
\renewcommand{\labelitemii}{$\circ$}



\title{}
\author{استاد درس:‌ دکتر جوانمردی، دکتر خادمیان \\ امیررضا کاظم لو \\ \small درس ربات های متحرکت خودگردان \\دانشگاه صنعتی مالک اشتر}
\date{\today}

\begin{document}
\maketitle
برای این که کالیبریشن انجام شود ما به chessboard نیاز داشتیم که از سایت https://calib.io/pages/camera-calibration-pattern-generator این صفحه تهیه شد که در فایل \textbf{calib.io_checker_200x150_8x11_15.pdf} قرار گرفته است و سپس آن را در صفحه نمایش کامپیوتر نمایش داده شده است و از زوایای مختلف و فاصله های مختلف تصویر برداری شده است. نکته مهم در تصویری که در زیر مشخص شده است:
\begin{figure}[H]
	source= "./clarification the size of each square.jpg"
\end{figure}
اندازه هر کدام از مربع ها برابر ۱۶ میلی متر هست و این مبنای ما در تبدیل سه بعد به دو بعد و محاسبه ماتریس درونی و بیرونی دوربین هست.
در بخشی از تمرین باید مقدار square_size_mm را مشخص می کردیم که این مقدار از قبل بر اساس متر مشخص شده بود ولی این مقدار می تواند بر اساس میلی متر هم مشخص شود تنها نکته ای که وجود دارد این است که اگر این متغیر را بر اساس متر مشخص کنیم خروجی ها نیز بر اساس متر می شود و اگر بر اساس میلی متر مشخص کنیم خروجی ها نیز بر اساس میلی متر می شود که در کد های تهیه شده توسط من به صورت متر دکر کرده ام یعنی ۰.۰۱۶.


\end{document}